\chapter{Discussion}

The studies which comprise this thesis leave its author with a several key insights into \textsc{ppi} inference methods that underscore their current limitations and paths towards addressing them. These insights are outlined here in detail.

\section*{Appropriate regularization methods}
\addcontentsline{toc}{section}{Appropriate regularization methods}

The application of regularization is a central component to the out-of-distribution performance of both \textsc{rapppid} and \textsc{intrepppid}, as demonstrated by ablation experiments. To be effective, however, the method of regularization must be chosen to appropriately reflect the architecture of the model and the underlying data. \textsc{rapppid} and \textsc{intrepppid} offer an interesting look into how progressively overlapping, specifically-targeted regularization can result in gains in model robustness and generalization. 

\subsection*{Regularization in RAPPPID}
\addcontentsline{toc}{subsection}{Regularization in \textsc{rapppid}}

Much of the regularization applied in the training of \textsc{rapppid} emerged from work by Stephen Merity \textit{et al.} on adapting existing methods of regularization for use with \textsc{lstm}s \cite{awd_lstm}. Among the methods Merity \textit{et al.} describe, the one that resulted in the greatest improvement in their testing is employing weight-dropped \textsc{lstm}s. Weight-dropping, also known as DropConnect, is a method of regularization similar to Dropout, but operates at the level of model weights rather than network activations \cite{dropconnect, dropout}. Comparisons to other architectures show that dropping out the hidden activations of \textsc{lstm}s greatly underperforms weight-dropping \cite{awd_lstm}. Indeed, applying dropout to \textsc{lstm} activations has been identified for quite some as a point of frustration when training recurrent models  time, resulting in ``dramatic'' changes in model dynamics \cite{fast_dropout}.

%Additionally, Merity \textit{et al.} rightly acknowledge the neutral performance impact. Randomly zeroing out weights rather than activations allows for the use of so-called ``black-box'' \textsc{lstm} implementations which have efficient low-level firmware libraries which are provided by hardware providers (such as \textsc{nvidia}'s cu\textsc{dnn}). 

In our own ablation experiments conducted as part of the development of \textsc{rapppid}, we found that tokenizing amino acid sequences had the greatest effect on model performance on an out-of-distribution testing set. Tokenization improves \textsc{rapppid}s performance by simultaneously reducing the number of timesteps, thereby mitigating the risk of vanishing gradients. Tokenization increases the pool of possible inputs at every timestep, as they no longer are limited to 22 amino acids, but any number of their permutations. Increasing the vocabulary like so can exacerbate the risk of overfitting. Luckily in the case of \textsc{rapppid}, the regularization opportunities afforded by tokenization greatly outweighed the risk of overfitting. The Unigram tokenization algorithm via the SentencePiece library was used to tokenize amino acid sequences. This algorithm allows for subword regularization, a process by which multiple segmentations for a sequence are sampled. For instance, the first 10 residues of the sequence for Human insulin (\texttt{MALWMRLLPL}) can be represented by the following sequence of tokens:
\begin{itemize}
    \item{\texttt{M}, \texttt{AL}, \texttt{W}, \texttt{M}, \texttt{R}, \texttt{L}, \texttt{L}, \texttt{PL}}
    \item{\texttt{M}, \texttt{A}, \texttt{L}, \texttt{W}, \texttt{M}, \texttt{R}, \texttt{LL}, \texttt{P}, \texttt{L}}
    \item{\texttt{M}, \texttt{A}, \texttt{L}, \texttt{W}, \texttt{M}, \texttt{RL}, \texttt{L}, \texttt{P}, \texttt{L}}
\end{itemize}
and \textit{etc}. One of all possible segmentations of the sequence are then chosen at random and used as the input to the model in question. The noise introduced by this process is regularizing and helps avoid the problem of rare tokens, whereby some tokens are under sampled at training time. Additionally, as part of the recommendations of the Merity \textit{et al.} study, ``embedding dropout'' is employed, which is a process by which the embeddings of certain tokens are randomly zeroed out. This regularizes the network in a similar way to random dropout on the activations or weights of a network, but particularly targets the problem of overfitting to the large vocabulary that was created using the Unigram algorithm.

\subsection*{Regularization in INTREPPPID}
\addcontentsline{toc}{subsection}{Regularization in \textsc{intrepppid}}

The focus on choosing regularization methods for \textsc{rapppid} centred on their appropriateness for the overall architecture. \textsc{intrepppid} inherited the architecture-focused regularization of \textsc{rapppid}, but also explored regularization methods that were informed by biological expert knowledge and inductive biases. Specifically, we sought to leverage our knowledge of orthologous protein sequences to improve the generalization of the model on testing examples from species other than those it was trained. 

We took inspiration from a common regularization technique used in the field of computer vision where images are subtly augmented such that many pixels are altered, but their semantic meaning is unchanged \cite{image_augmentation}. Orthologous proteins are similarly altered subtly by evolution, but typically retain their function and purpose. Incidentally, the application of augmentation techniques to proteins was simultaneously explored by Sun \textit{et al.}, albeit by directly and randomly altering sequences rather than by using orthologous proteins \cite{protein_augmentation}. This direct mutation approach scales much more easily than our orthologous method, as you are not limited by the number of orthologous sequences found in curated knowledge-bases. Unlike using orthologous protein sequences, randomly altering sequences can however result in proteins which lose their semantic meaning --- much as an image of the number 9, when rotated $180^{\circ}$, loses its semantic meaning and becomes a 6.

The final result of adopting the orthologous locality loss introduced in the \textsc{intrepppid} study is a regularized latent space where proteins which are similar in function across species share similar vector embeddings.

\section*{Questions regarding sequence identity and its role in data leakage}
\addcontentsline{toc}{section}{Questions regarding sequence identity and its role in data leakage}

The three studies reproduced in this thesis are greatly indebted to the work of Park and Marcotte, which identified their strictest dataset type (\textsc{c3}) as having ``test pairs sharing neither protein with the training set'' \cite{park_marcotte}. Membership of a protein in either the testing, validation, set can be defined a number of ways. For instance, membership can be determined by shared UniprotKB identifiers of proteins. The trouble, however, in using protein names and most identifiers is that many protein names, while unique, have sequences that are identical, or nearly so. For instance, the January 2025 release of Uniprot has two entries for the human protein ``Cytochrome c oxidase subunit 3'' with identifiers \texttt{A0A023Q4K2} and \texttt{Q4EZD1} which differ by one additional serine residue in the former. 

To help identify proteins that are identical in all but name, clustered identifiers such as those pre-computed in the UnirRf100 database, are commonly used. The two Uniprot identifiers for human Cytochrome c oxidase subunit 3 map to a single UniRef100 identifier \\(\texttt{UniRef100\_A0A023Q4K2}). By mapping Uniprot identifiers to UniRef100 identifiers, it is possible to avoid considering proteins with identical sequences as distinct proteins. 

A consequence of the high degree of homology between proteins is that many proteins approach, but fall short of 100\% sequence identity. These homologous proteins often share the same function and interact with similar proteins, and therefore likely are a source of data leakage and warrant removal. Highly-conserved proteins like insulin (\texttt{P01308}) are prolific examples of this phenomena, sharing at least 90\% sequence identity with 14 other proteins in Uniprot. Matters get much worse when one expands the similarity threshold to  at least 50\%, revealing 53 other homologous proteins in Uniprot from species as diverse as \textit{Rattus norvegicus} (Rat) and \textit{Gorilla gorilla gorilla} (Western lowland gorilla). Once more, the UniRef database allows us to map similar proteins to the same identifiers. For example, proteins with at least 90\% sequence identity and 50\% sequence identity to insulin can be mapped to the identifiers \texttt{UniRef90\_P01308} and \texttt{UniRef50\_P01308}, respectively.

Park and Marcotte used, without a stated justification, a threshold of 40\% sequence identity to determine protein similarity \cite{park_marcotte}. This threshold is used by many others when establishing a similarly strict training/testing split without much or often any discussion as to how it was selected. For instance, Bernett \textit{et al.} used the \textsc{cd-hit} clustering software to establish validation folds with no more than 40\% sequence identity between them \cite{black_box, cd_hit}. Prior to that study, Hamp and Rost used a sequence-structure alignment database (\textsc{hssp}) ``...correspond[ing] to 40\% pairwise sequence identity for 250 aligned residues" \cite{Hamp_Rost_2015, Sander_Schneider_1991, hssp}.

While the rationale for these seemingly arbitrary thresholds is not provided in those studies, one begins to emerge after a thorough review of the relevant literature, though there is no guarantee it is shared with the original authors. It is possible that these thresholds, particularly in the case of Hamp and Rost, were influenced by a 1991 study by Sander and Schneider which found that protein sequences with more than 30\% sequence identity often share similar structure \cite{Sander_Schneider_1991}. Subsequent work by Burkhard Rost found that proteins identified as having between 20-35\% sequence identity, a so-called protein-sequence alignment ``twilight zone'', are far more likely to be ``false negatives'' \cite{rost_twilight_zone}. In this context, false negatives are taken to mean proteins which, despite falling below the sequence identity threshold, are structurally similar and false positive are proteins above the threshold which do not share structural similarity. Rost reports:

\begin{displayquote}
\begin{singlespace}
The number of false positives rose by a factor of 5 000 (Figure 2B). Hence, separating true and false positives switched from a trivial task (above 35\% [pairwise sequence identity]) to the problem of finding needles in a haystack (20–30\%). \cite{rost_twilight_zone}
\end{singlespace}
\end{displayquote}

In light of these findings, a threshold of 40\% sequence identity appears to be a fair compromise between identifying as many structurally similar proteins as possible while greatly reducing the number of proteins false-positives.

While we have identified a possible rationale for why the 40\% sequence identity threshold has been identified --- it is a proxy for the presence of structural similarity --- we have yet to examine any evidence as to whether it is reasonable or useful. I argue here that eliminating nonindependence between the training and testing set must be balanced with creating unduly biased datasets.

Supervised machine learning algorithms rely on example input-output pairs to learn a generalized mapping between the two. These training examples are required to be representative of the variation observed in the universe of input-output pairs that are expected to be correctly inferred. To understand whether a trained model is performing at the level intended, a set of evaluation input-output pairs must be chosen. If it is to give us an accurate representation of the model's performance, the evaluation set must simultaneously (1) be representative of the problem space and (2) avoid data leakage. Data leakage is a topic well discussed throughout this thesis, however it is useful here to use the taxonomy outlined by Kapoor and Narayanan to identify some of its common subtypes \cite{Kapoor_Narayanan_2023}. In the framework of this taxonomy, the Park and Marcotte study identified examples of  \textsc{l3.2} and \textsc{l1.4} errors --- \textit{i.e.}, ``nonindependence between training and test samples'' and ``duplicates in datasets'', respectively \cite{park_marcotte, Kapoor_Narayanan_2023}.

The \textsc{l1.4}, or duplication, error describes instances where duplicate data is present in the dataset increasing the risk that ``the same data could exist in the training set and the test set'' \cite{Kapoor_Narayanan_2023}. As explored previously, homologous sequences result in nearly identical sequences, and therefore \textsc{ppi} datasets which do not account for homology will almost certainly observe an \textsc{l1.4} data leakage. By applying a maximum sequence identity threshold between testing and training datasets of 40\%, Park and Marcotte side-step the issue of duplicates. A 100\% threshold, however, would also eliminate the presence of duplicate sequences. A further 5-10\% allowance for small variations corresponding to threshold between 95-90\% would be appropriately sufficient to eliminate near-duplicates.

The \textsc{l3.2}, or nonindependence, error describes instances where testing samples are not drawn independently from the distribution of the problem space, but rather from one present in the training set that is not central to the scientific claim of the model. Examples provided by Kapoor and Narayanan include histology classification models which are evaluated on datasets of patients also present in the training set, or protein function classification models where a given protein family is present in both training and testing sets \cite{Oner_Cheng_Lee_Sung_2020, Whalen_Schreiber_Noble_Pollard_2022}. In the case of \textsc{ppi} inference models, it is uncontroversial to assert that Park and Marcotte identified the \textsc{l3.2} error caused by shared proteins in testing and training sets and remediated it by removing the dependence between the training and testing examples. It remains to be explored, however, which precise dependency was eliminated. Was it a dependency between shared \textit{proteins} in the testing and training set, or rather a dependency between \textit{homologous} proteins? Depending on one's interpretation of the role of the chosen sequence identity threshold, perhaps even a dependency on shared protein families is being avoided. If the former, most threshold values would be sufficient to mitigate the error. If the latter, then 40\% would appear to be somewhat of an upper-bound among threshold values that would capture structural similarity, as was discussed earlier in reviewing the works of Sander and Schneider and Rost \cite{Sander_Schneider_1991, rost_twilight_zone}.
\\\\\\
The murkiness of this question is to be anticipated, according to Kapoor and Narayanan, who warn:

\begin{displayquote}
\begin{singlespace}
Handling nonindependence between the training and test sets in general (\textit{i.e.}, without any assumptions about independence in the data) is a hard problem, because we might not know the underlying dependency structure of the task in many cases.\cite{Kapoor_Narayanan_2023}
\end{singlespace}
\end{displayquote}

Indeed, it is not easy to disentangle which aspects of the training and evaluation set distribution is intrinsic to learning a generalized model, and which serve only to artificially inflate the testing performance of the model via data leakage. While it may be tempting at first to resort to lower thresholds in the range of 30-50\%, as many studies have done, there is a risk that models are being needlessly deprived of training data causing gaps in the coverage protein families which have representative samples in the training set. An over-correction could result in exacerbating the risk of biasing the training data, reducing its magnitude and variance, and ultimately resulting in overfitting.

Additionally, there is a question of whether stringent sequence identity thresholds are effective methods for excluding proteins on the basis of their protein family membership. Rost reminds us that his findings echo those of the earlier Sander and Schneider study: ``...high levels of sequence similarity or identity do \textit{not} ascertain structural similarity...'' \cite{rost_twilight_zone}.

Further study into the extent of dependence between the training and testing splits of \textsc{ppi} datasets, at what level it occurs, and possible mitigation strategies, is gravely needed. Until then, interpretation of the findings in the works presented herein offer a suggestion of what such a study might confirm. 

In the study in which it was introduced, \textsc{rapppid} was trained and tested on a \textsc{c3}-style dataset, however protein homology was not accounted for \cite{rapppid}. This is to say, proteins were excluded solely on the basis of shared \textsc{string} protein identifiers across training and testing datasets. The \textsc{intrepppid} study, by contrast, saw a \textsc{rapppid} model trained on a \textsc{c3}-style dataset with a 90\% maximum sequence identity threshold between training and testing splits \cite{intrepppid}. Introducing the 90\% maximum sequence identity threshold results in \textsc{ppi} inference performance, as measured by \textsc{auroc}, dropping by $\approx8\%$. The last study presented in this thesis includes an analysis in the change of performance of p\textsc{llm} models when trained on \textsc{ppi} datasets with maximum sequence identity threshold between their training and testing sets ranging between 50\% to 90\% \cite{szymborski_emad_2025}. Over that range, very little variation is observed. While further studies are required to draw any conclusions, we are left with the suggestion that, while it is pertinent to account for protein homology, little if any data leakage is mitigated by varying the maximum sequence identity threshold over the range of 50\% to 90\%.

\section*{Distal impacts of efficient architectures}
\addcontentsline{toc}{section}{Distal impacts of efficient architectures}

\subsection*{As a means to improve generalization}
\addcontentsline{toc}{subsection}{As a means to improve generalization}

\textsc{rapppid}'s encoder, which make up the majority of the network by neuron count, is constructed using a twin neural network architecture. By sharing weights between the encoders of each protein, \textsc{rapppid} effectively halves the number of weights needed to infer \textsc{ppi}s. While the forward pass requires the same number of floating point operations (\textsc{flops}), the number of \textsc{flops} required for the backward-pass is halved. Seeing as the back-propagation is far more computational-intensive than the forward pass, the backward-pass is the more relevant target for optimization. A study by Richoux \textit{et al.} confirms that twin neural network constructions observe far less overfitting than a twin construction \cite{richoux}. This finding exemplifies the so-called ``bias/variance dilemma'' that relates the increase in neurons of a neural network to an increase in variance, leading to overfitting \cite{bias_variance}.

The \textsc{intrepppid} expands this concept by adding an additional three passes to compute an orthology-related task, forming what we've come to call a ``quintuplet'' network. This allows us to leave the number of effective weights in the network unchanged, increase the amount of data the network can learn from, and add an additional source of inductive bias. The result is a network that generalizes better than \textsc{rapppid}, despite maintaining the same number of weights, requires the same number of \textsc{flops} for \textsc{ppi} inference, and requires a comparable amount of time to train.

High-quality, unbiased sources of biological data are difficult to come by due to the cost associated with generating such data, and \textsc{ppi} pairs are no exception. After filtering out low confidence interaction pairs and creating a strict validation dataset, only 39 229 training \textsc{ppi} pairs were used to train \textsc{intrepppid}. This stands in comparison to the 220 000 parameters that make-up \textsc{intrepppid}'s encoder, alone. By increasing the amount of data the model is able to learn from (in the form of the orthologous locality task) while leaving the number of parameters unchanged, \textsc{intrepppid} reduces the degree of over-parameterization that is common among contemporaneous \textsc{ppi} inference methods.

It's important to note, however, that recent studies have found that over-parameterization in deep neural networks, past some (ill-defined) point, can cease to cause overfitting and result in well-calibrated, generalized models \cite{Baldi_2024}. Further, it has been found that the variance-bias trade-off is not universally applicable \cite{bias_variance_update}. These developments are, in part, leveraged in the study presented in Chapter 5. Briefly, p\textsc{llm}'s offer an interesting solution to the over-parameterization problem that plagues \textsc{ppi} inference models, along with many other protein-related tasks which suffer from constrained training datasets.


%\subsection*{Headroom for Expansion}

\subsection*{Lowering barriers to PPI inference research}
\addcontentsline{toc}{subsection}{Lowering barriers to \textsc{ppi} inference research}

\begin{displayquote}
\begin{singlespace}
\textit{``Necessity is the mother of invention.''}
\\\hspace*{6mm}--- Proverb
\end{singlespace}
\end{displayquote}

Computational resources are not equally distributed amongst researchers, not at the global scale nor most regional scales. An indication of this is that the top 10 super-computers by performance on the \textsc{linpack} benchmark all belong to countries in the Global North, according to \textsc{top500} as of June 2025 \cite{TOP500_2025, linpack}. While not a perfect measure of the accessibility of computational resources, it follows that access is greater in wealthy countries that can and do allocate more funds to public and private research in the form of grants and tax credits.

Access to computational hardware is an important factor in being able to explore and successfully study research questions in the field of deep learning. In her influential essay ``The Hardware Lottery'', Sara Hooker argues that ``[i]n the field of artificial intelligence research ... it is our tooling which has played a disproportionate role in deciding
what ideas succeed (and which fail)'' \cite{hardware_lottery}. Over the course of my doctoral studies, I've had the good fortune to have had far greater access to computational resources than the majority of individuals globally. Despite this, I had access to many orders of magnitude fewer \textsc{gpu} hours and video random-access memory (\textsc{vram}) than leading research groups, particularly in the private sector. To narrow this gap, I applied a great focus on efficient deep learning architectures to make the most of the resources made available to me. This approach is advocated by Hooker as a way to mitigate the disparity in access to computational resources:

\begin{displayquote}
\begin{singlespace}
Software needs to do more work, but it is also well positioned to do so. We have neglected efficient software throughout the era of Moore’s law, trusting that predictable gains in compute would compensate for inefficiencies in the software stack. This means there are many low hanging fruit as we begin to optimize for more efficient code... \cite{hardware_lottery}
\end{singlespace}
\end{displayquote}

A out-sized contributor to this hardware inequity on the side of software is the overwhelming popularity of the compute-hungry transformer architecture \cite{aiayn}. An oft-cited concern regarding the computational intensity of transformers is the self-attention mechanism at their heart, whose complexity scales quadratically with the length of its input sequence \cite{long_range_arena, Zimerman_Wolf_2023}. As a result, the attention mechanism of the transformer architecture is often the target of improvement and optimization \cite{longnet, flashattention, linformer, longformer}. 

Despite this, Iandola \textit{et al.} found that the fully-connected layers in the feedforward network is responsible for the overwhelming share of \textsc{flops} (73\%) and latency on mobile devices (69.4\%) in traditional transformer architectures \cite{squeezebert}. They proceed to develop a convolution-based alternative to the fully-connected layers in the feedforward modules of the transformer architecture \cite{squeezebert}. Using this construction, we were capable of training SqueezeProt, a p\textsc{llm} with less than 7\% of the number of parameters of  leading p\textsc{llm} Prot\textsc{bert}, computes tokens $5\times$ faster, and achieves comparable \textsc{ppi} inference accuracy. \textsc{vram} is often the primary hardware bottleneck when training large neural networks and is commonly the main price differentiator between \textsc{gpu}s of the same hardware architecture. While SqueezeProt-\textsc{u50} was trained on four Nvidia \textsc{rtx} 3090 \textsc{gpus}, the relatively-low number of parameters allows for SqueezeProt to be trained on one such \textsc{gpu}. The \textsc{rtx} 3090 is a consumer \textsc{gpu} which had an \textsc{msrp} of \$1 499 \textsc{usd} at launch \cite{rtx_3090_msrp}. By comparison, Prot\textsc{bert} was trained on two supercomputers --- one of which is, as its authors describe, ``the world’s second fastest computer'' comprised of 27 708 Nvidia Tesla \textsc{v100} \textsc{gpu}s \cite{elnaggar_prottrans_2022}. The Nvidia Tesla \textsc{v100} is a data center \textsc{gpu} with an \textsc{msrp} of \$10 000 \textsc{usd} at launch \cite{Keefer_2025}.

Through the early adoption of efficient neural network architectures, we were able to conduct research on topics which are dominated by laboratories with asymmetric access to computational resources. Despite this victory, the inequities of the hardware lottery are unlikely to be optimized away by efficient architectures like the ones presented in this thesis. Hooker concludes:

\begin{displayquote}
\begin{singlespace}
...the hardware lottery has not gone away, and the gap between the winners and losers will grow increasingly larger. In order to avoid future hardware lotteries, we need to make it easier to quantify the opportunity cost of settling for the hardware and software we have. \cite{hardware_lottery}
\end{singlespace}
\end{displayquote}

\subsection*{Environmental impacts}
\addcontentsline{toc}{subsection}{Environmental impacts}

\begin{displayquote}
\begin{singlespace}
\textit{``...invention is the mother of necessity.''}
\\\hspace*{6mm}--- Thorstein Veblen\\\hspace*{11mm}The Instinct of Workmanship and the State of the Industrial Arts
\end{singlespace}
\end{displayquote}

% To produce a two-gram memory microchip, 1.3 kilograms of fossil fuels and materials are required.

%While the substances used in the manufacture of microchips contribute to pollution at both bodily and environmental levels, this condition intensifies with the sheer quantity of microchips manufactured

As established in the previous subsection, deep learning is a power-hungry endeavour that relies on specialized hardware. Both the high energy and semiconductor demands of deep learning algorithms have associated environmental footprint. In this subsection, I explore briefly what the relevance of the efficient architectures adopted in this thesis may have on their relative environmental impact.

On the topic of energy consumption, both SqueezeProt and the top-ranking p\textsc{llm}s analysed herein cause an Nvidia \textsc{rtx} 3090 \textsc{gpu} to consume just under its Thermal Design Power (\textsc{tdp}) of 350 watts. SqueezeProt can compute embeddings at a far faster rate than the top-ranking p\textsc{llm}s, however, which results in energy requirements for embedding computation that is orders of magnitude lower for the same amino acid sequence. Specifically, SqueezeProt requires approximately 8.33 millijoules (mJ) per amino acid embedded whereas Prot\textsc{bert}, \textsc{esm2}, and Prot\textsc{t5} require 37.4 mJ, 52.3 mJ, and 67.5 mJ, respectively. As a practical example, it would require the same energy to power the average Canadian home for approximately 4 days to compute embeddings for all 253 061 696 proteins currently in UniProt using SqueezeProt\footnote{Assuming an average amino acid length of 400 amino acids and a yearly household energy consumption of 85.4 gigajoules \cite{the_uniprot_consortium_uniprot_2023, statscan_energy_consumption}.}. This number balloons to over two weeks (16 days) in the case of Prot\textsc{bert} and just under a month (29 days) for Prot\textsc{t5}. By eschewing transformers entirely, power consumption could be greatly reduced. \textsc{rapppid} causes an \textsc{rtx} 3090 to consume about half of its \textsc{tdp} according to my measurements, resulting in the use of a mere 0.543 mJ per amino acid embedded. In our UniProt-wide embedding scenario, \textsc{rapppid} would require the same energy as the typical Canadian household uses in 5 hours and 36 minutes. Notably, all network training in this study was conducted in the province of Qu\'ebec where 99\% of power delivered to the energy grid is produced by non-emitting sources \cite{canada_energy_regulator_canadas_2021}.

Energy consumption is but one of the environmental consequences of deep-learning-based \textsc{ppi} inference methods. Another important ecological concern is the life-cycle impact of the manufacturing of the specialized semiconductors required to train and operate these models. According to Jennifer Gabrys, Chair in Media, Culture and Environment at the University of Cambridge, ``[t]o produce a two-gram memory microchip, 1.3 kilograms of fossil fuels and materials are required'' \cite{silicon_elephants}. The aforementioned materials are severely injurious to human and animal health, with sweeping implications for individuals born in regions contaminated by electronic waste --- including, but not limited to: ``...suppressed immune response, hearing loss, altered hepatic and renal function, decline in oral anti-inflammatory ability, and abnormal [fasting blood glucose level]'' as well as adverse birth outcomes \cite{Parvez2021}. By reducing the dependence on such specialized semiconductor hardware, we can hope to lessen the ecological and health impacts caused by their manufacturing and disposal. Gabrys writes:

\begin{displayquote}
\begin{singlespace}
While the substances used in the manufacture of microchips contribute to pollution at both bodily and environmental levels, this condition intensifies with the sheer quantity of microchips manufactured. \cite{silicon_elephants}
\end{singlespace}
\end{displayquote}

As discussed in the previous sub-section, \textsc{vram} is often the most relevant bottleneck when considering \textsc{gpu}s. Models, particularly \textsc{llm}s, often require more \textsc{vram} than afforded by one \textsc{gpu}. For this reason, multiple \textsc{gpu}s are often required. By adopting an efficient architecture which requires a small fraction of the number of parameters of other leading p\textsc{llm}s, it is possible to train SqueezeProt on a single consumer-class \textsc{gpu} rather than many.

%https://www.semanticscholar.org/paper/Recycling-of-Laboratory-Plastic-Waste-%E2%80%93-a-Study-on-Martin-Sover/691d87638cf4401dff19009506c12ec52be48e19

%https://www.semanticscholar.org/paper/Plastic-Pollution-and-Sustainable-Managing-of-Waste-Tan-Ying/61cbfcf93839c00797b1b40915c36778e64d231c

%https://www.semanticscholar.org/paper/A-case-report%3A-insights-into-reducing-plastic-waste-Alves-Sargison/8d35d3710814bffa0c9e198b5bdee57ac93b6b0d

%https://www.semanticscholar.org/paper/Medical-laboratory-waste-generation-rate%2C-practices-Endris-Tamir/98f22ec71f43d2647463926b6f0ddaebc121f0e4

%https://www.semanticscholar.org/paper/Approaching-sustainability-in-Laboratory-Medicine-Thakur-Mukhopadhyay/67c63e90071051c98758fa8fbb38642efc5f621b

It may be tempting to see these economies in the amount of energy and number of microchips required to train and operate the models described in this thesis and declare that there has been an overall improvement in the environmental impact of the model --- that the ecological consequences have been partly mitigated. To do so would be premature. Merely observing gains of efficiency is not sufficient to come to such a conclusion. Gains in efficiency of a process can indeed increase the demand and usage of its inputs. William S. Jevons famously observed that the demand, and consequently the usage, of iron and coal increased as machines who required them became more efficient. Jevons' paradox, as it so became known, can be imagined in the context of this work: if p\textsc{llm}-based \textsc{ppi} inference methods become more accurate and less costly, researchers and industry may adopt them where they had not before. This would result in energy and environment costs that would not have been incurred if p\textsc{llm}-based \textsc{ppi} inference methods were inefficient, and therefore less desirable. Jevons' paradox, and the so called ``rebound'' debate which has arisen from it, is far from a settled matter and lacks broad expert consensus \cite{alcott_jevons_2005}. Also to be considered are the environmental impacts of the waste and pollution caused by the experimental techniques used to discover \textsc{ppi}s that inference methods might replace \cite{endris_medical_2022, alves_case_2021, thakur_approaching_2024, tan_plastic_2022}.

Whether p\textsc{llm}-based \textsc{ppi} inference, even efficient ones as I've described, are an improvement over their alternatives requires systematic study. Frameworks, such as Chistopher Becker's five factors of just sustainability design (\textsc{jsd}), are a useful basis for such a study \cite{becker_what_2023}. Regrettably, performing such a thorough analysis could demand another doctoral degree's worth of study, which I have deemed out of scope.